\documentclass[12pt, a4paper]{article}

\usepackage[margin=2.5cm]{geometry}
\usepackage{graphicx}
\usepackage{hyperref}
\usepackage{parskip}
\usepackage{setspace}
\usepackage{titlesec}
\usepackage{booktabs}
\usepackage[round]{natbib}
\usepackage{microtype}
\usepackage{xcolor}

\hypersetup{
  colorlinks=true,
  linkcolor=blue!60!black,
  citecolor=blue!60!black,
  urlcolor=blue!60!black
}

\titleformat{\section}{\large\bfseries}{}{0em}{}[\titlerule]
\titleformat{\subsection}{\normalsize\bfseries}{}{0em}{}

\setstretch{1.2}

\title{\textbf{Week 1 Report}\\[0.4em]
  \large Algorithmic Trading --- Summer of Science}
\author{Arush Anand}
\date{May 2026}

\begin{document}

\maketitle
\tableofcontents
\newpage

%% ──────────────────────────────────────────────────────────────────
\section{Chapter 1 Summary: \textit{Quantitative Trading} by E.\ P.\ Chan}
%% ──────────────────────────────────────────────────────────────────

\subsection{What Is Quantitative Trading?}

\citet{chan2021} says that Quantitative Trading is the trading of 
securities based strictly on the buy/sell decisions of computer 
algorithms. It is not quite a technical analysis as only metrics
that are quantifiable (as opposed to them being qualitative) are
used in Quantitative Trading. 


\subsection{Who Can Become a Quantitative Trader?}

In his own words, 
"The ideal independent quantitative trader is therefore someone
who has some prior experience with finance or computer programming,
who has enough savings to withstand the inevitable losses and
periods without income, and whose emotion has found the right balance 
between fear and greed." This ideal trader might not have the academic
qualifications of the stereotypical algorithmic trader, as thats not
a prerequisite or even that big of an edge over the layman.

\subsection{The Business Case for Quantitative Trading}

Chan frames a quantitative trading operation as a small business and
then identifies three ways in which it differs favourably from most
other ventures.

\subsubsection*{Scalability}
Increasing the size of a quantitative portfolio often requires nothing
more than adjusting a leverage parameter in code. No venture capital,
no bank negotiation. Proprietary trading firms can offer intraday
leverage of up to $\times$40 on equity, while futures and currency
accounts routinely provide $\times$10 or more from standard brokerages.
Chan warns, however, that leverage is a double-edged sword and should
never be chased for its own sake.

\subsubsection*{Demand on Time}
Automation is intrinsic to algorithmic trading. In Chan's own workflow,
day-to-day operations (data downloads, signal generation, order
transmission) initially consumed roughly two and a half hours around
the market open and close; over time these were fully automated until
active intervention fell to zero. Research and strategy development 
can be scheduled flexibly. The largest psychological challenge, 
Chan admits, is resisting the urge to manually override positions 
during intraday drawdowns. 

\subsubsection*{No Marketing Required}
In most businesses revenue depends on persuading customers. In trading,
the counterparty's only concern is price. As long as a trader is not
managing outside money, the product (the strategy) markets itself
through its P\&L. This freedom from salesmanship is, for many
practitioners, the greatest appeal of the profession.

\bigskip
\noindent\textbf{Key takeaways from Chapter~1:}
\begin{itemize}
  \item Quantitative trading is a form of systematic, algorithm-driven trading.
  \item In a search for complex strategies, it is usually the simpler ones which would benefit us more.
  \item The business advantages are scalability, time efficiency through automation, and zero marketing overhead.
\end{itemize}

%% ──────────────────────────────────────────────────────────────────
\newpage
\section{Kaggle Pandas Certificate}
%% ──────────────────────────────────────────────────────────────────

As part of the Week~1 curriculum I completed the \textbf{Pandas} course
on Kaggle, which covers the core data-manipulation library used
throughout quantitative research in Python. The certificate of
completion is reproduced below.

\begin{figure}[h]
  \centering
  \includegraphics[width=\textwidth]{{Arush Anand - Pandas}}
  \caption{Kaggle Pandas course certificate --- Arush Anand.}
  \label{fig:certificate}
\end{figure}

%% ──────────────────────────────────────────────────────────────────
\bibliographystyle{plainnat}
\bibliography{references}

\end{document}
